%% Separately authored historical derivative following exact source lines 10353--10674.
%% Necessity: exact restoration replaces the historical ordinal, hierarchy, and
%% collapse rewrites, but it does not contain the separately developed ordinal
%% arithmetic, stronger rank consequences, or hierarchy diagram retained here.
%% Definition: this file is based on historical body slices 288--460, 505--514,
%% and 530--564 in that order; it adds the descriptive section boundary and
%% removes internal apparatus IDs from one table and two figure captions.
%% Standing: this is separately authored reader material, not inherited source
%% or a proof-status upgrade.
%% Nonredundancy: the exact source supplies the hierarchy definitions and proofs;
%% these slices instead supply ordinal arithmetic, an additional rank theorem,
%% and apparatus A-129/A-130.
%% Use: parts/p06.tex routes this file once after the four exact chapter wrappers.

% BEGIN HISTORICAL SLICE 288-460
\section{Ordinal arithmetic and Cantor normal form}
\label{sec:p06-ordinal-arithmetic}

\begin{hbdefinition}{def:p06-ordinal-arithmetic}{Ordinal addition, multiplication, and exponentiation}
Ordinal arithmetic is defined by transfinite recursion on the \emph{right}
argument.  For a nonzero limit ordinal \(\lambda\),
\[
\begin{array}{lll}
 \alpha+0=\alpha,
 &\alpha+(\beta+1)=(\alpha+\beta)+1,
 &\alpha+\lambda=\sup_{\xi<\lambda}(\alpha+\xi),\\[2pt]
 \alpha\cdot0=0,
 &\alpha\cdot(\beta+1)=\alpha\cdot\beta+\alpha,
 &\alpha\cdot\lambda=\sup_{\xi<\lambda}(\alpha\cdot\xi).
\end{array}
\]
For \(\alpha>0\), exponentiation is
\[
 \alpha^0=1,
 \qquad
 \alpha^{\beta+1}=\alpha^\beta\cdot\alpha,
 \qquad
 \alpha^\lambda=\sup_{\xi<\lambda}\alpha^\xi;
\]
and \(0^0=1\), while \(0^\beta=0\) for \(\beta>0\).  These are ordinal
operations; they are not cardinal addition, multiplication, or exponentiation.
\index{ordinal addition}\index{ordinal multiplication}
\index{ordinal exponentiation}
\end{hbdefinition}

\begin{hbtheorem}{thm:p06-ordinal-arithmetic-laws}{cited=Jech2003}{Basic laws of ordinal arithmetic}
For ordinals \(\alpha,\beta,\gamma\), addition and multiplication are
associative,
\[
 \alpha\cdot(\beta+\gamma)=
 \alpha\cdot\beta+\alpha\cdot\gamma,
 \qquad
 \alpha^{\beta+\gamma}=\alpha^\beta\cdot\alpha^\gamma.
\]
For fixed left argument, all three operations satisfy their defining supremum
clause at nonzero limits (for exponentiation, with nonzero base).  Addition and
multiplication are not commutative, and multiplication need not distribute in
the other direction.  The exact law-and-failure register is
Table~\ref{tab:p06-ordinal-arithmetic}.
\end{hbtheorem}

\begingroup\small
\begin{longtable}{@{}p{0.19\textwidth}p{0.46\textwidth}p{0.27\textwidth}@{}}
\caption{Ordinal-arithmetic laws, continuity clauses, and diagnostic
failures.}\label{tab:p06-ordinal-arithmetic}\\
\toprule
Operation & Valid laws and limit behavior & Failures and cautions \\
\midrule
\endfirsthead
\caption[]{Ordinal-arithmetic laws, continuity clauses, and diagnostic
failures (continued)}\\
\toprule
Operation & Valid laws and limit behavior & Failures and cautions \\
\midrule
\endhead
\midrule\multicolumn{3}{r}{\emph{continued on next page}}\\
\endfoot
\endlastfoot
Addition &
\(0+\alpha=\alpha=\alpha+0\); associativity;
\(\beta<\gamma\) implies \(\alpha+\beta<\alpha+\gamma\);
\(\alpha+\lambda=\sup_{\xi<\lambda}(\alpha+\xi)\). &
Not commutative: \(1+\omega=\omega<\omega+1\). \\
\addlinespace
Multiplication &
\(1\cdot\alpha=\alpha=\alpha\cdot1\); associativity;
\(\alpha(\beta+\gamma)=\alpha\beta+\alpha\gamma\).
If \(\alpha>0\), then
\(\beta<\gamma\) implies \(\alpha\beta<\alpha\gamma\);
\(\alpha\lambda=\sup_{\xi<\lambda}\alpha\xi\). &
Not commutative: \(2\cdot\omega=\omega<\omega\cdot2\).
The reverse distributive law fails:
\((1+1)\omega=\omega\ne\omega+\omega\). \\
\addlinespace
Exponentiation &
For \(\alpha>0\),
\(\alpha^{\beta+\gamma}=\alpha^\beta\alpha^\gamma\) and
\((\alpha^\beta)^\gamma=\alpha^{\beta\gamma}\).
For \(\alpha\geq2\), exponentiation is strictly increasing in the exponent;
for \(\alpha>0\),
\(\alpha^\lambda=\sup_{\xi<\lambda}\alpha^\xi\). &
The base-zero clauses are exceptional.  Cardinal-exponent laws cannot be
substituted for ordinal-exponent laws. \\
\addlinespace
Cantor normal form &
Finite decreasing powers of \(\omega\), with positive finite coefficients on
the right.  Leading exponents expose the comparison and limit structure. &
The order of summands is essential; rearrangement generally changes the
ordinal. \\
\bottomrule
\end{longtable}
\endgroup

\begin{hbtheorem}{thm:p06-cantor-normal-form}{cited=Jech2003}{Cantor normal form}
Every nonzero ordinal \(\alpha\) has a unique expression
\[
 \alpha=
 \omega^{\beta_0}\cdot n_0+
 \omega^{\beta_1}\cdot n_1+\cdots+
 \omega^{\beta_{k-1}}\cdot n_{k-1},
\]
where \(k\geq1\) is finite,
\(\beta_0>\beta_1>\cdots>\beta_{k-1}\), and each
\(n_i\in\N\setminus\{0\}\).  The coefficient is written on the right because
ordinal multiplication is not commutative.  The empty sum represents zero.
This is the \emph{Cantor normal form} of \(\alpha\)
\cite{Cantor1897,Sierpinski1965}.
\index{Cantor normal form}
\end{hbtheorem}

\begin{hbexample}{def:p06-cnf-arithmetic}{Order sensitivity in Cantor normal form}
The leading \(\omega^2\)-term of the right summand absorbs lower-degree terms
on its left:
\[
 (\omega^3\cdot2+\omega\cdot5+7)
 +(\omega^2\cdot4+3)
 =\omega^3\cdot2+\omega^2\cdot4+3.
\]
Also \((\omega+1)\cdot\omega=\omega^2\), whereas
\(\omega\cdot(\omega+1)=\omega^2+\omega\).  These are computations with the
recursive clauses, not cardinal-arithmetic identities.
\end{hbexample}

\begin{hbdefinition}{def:p06-epsilon-zero}{The ordinal epsilon-zero}
An \emph{epsilon ordinal} is an ordinal \(\varepsilon\) satisfying
\(\omega^\varepsilon=\varepsilon\).  The least such ordinal is
\(\varepsilon_0\).  If \(\alpha_0=1\) and
\(\alpha_{n+1}=\omega^{\alpha_n}\), then
\[
 \varepsilon_0=\sup_{n<\omega}\alpha_n.
\]
Every ordinal below \(\varepsilon_0\) has a finite term obtained from Cantor
normal forms with recursively nested exponents, but there is no single finite
tower height sufficient for all of them.  The term names an ordinal fixed
point, not a proof-theoretic notation system; later proof-theory entries state
separately how notation systems represent ordinals below it.
\index{epsilon ordinal}\index{epsilon-zero@\(\varepsilon_0\)}
\end{hbdefinition}

\begin{figure}[htbp]
\centering
\begin{tikzpicture}[
  x=1cm,y=1cm,
  milestone/.style={draw,rounded corners,fill=black!4,inner sep=2.5pt,
                    align=center,font=\scriptsize},
  note/.style={align=center,font=\scriptsize}
]
  \draw[-{Latex[length=2.5mm]},thick] (0,0) -- (11.2,0);
  \foreach \x in {0.4,1.65,2.8,4.0,5.2,6.5,8.05,9.45,10.8}
    \draw (\x,0.10) -- (\x,-0.10);
  \node[milestone,above=3pt] at (0.4,0) {finite\\ordinals};
  \node[milestone,above=3pt] at (1.65,0) {$\omega$};
  \node[milestone,above=3pt] at (2.8,0) {$\omega+1$};
  \node[milestone,above=3pt] at (4.0,0) {$\omega\cdot2$};
  \node[milestone,above=3pt] at (5.2,0) {$\omega^2$};
  \node[milestone,above=3pt] at (6.5,0) {$\omega^\omega$};
  \node[milestone,above=3pt] at (8.05,0)
    {$\omega^{\omega^\omega}$};
  \node[milestone,above=3pt] at (9.45,0) {$\cdots$};
  \node[milestone,above=3pt] at (10.8,0) {$\varepsilon_0$};
  \node[note,below=5pt] at (4.2,0)
    {finite Cantor sums with exponents below the next milestone};
  \node[note,below=18pt] at (9.1,0)
    {finite towers of unbounded height};
  \node[note,below=5pt] at (9.7,0)
    {$\sup(1,\omega,\omega^\omega,\omega^{\omega^\omega},\ldots)$};
\end{tikzpicture}
\caption{Ordinal milestones below \(\varepsilon_0\) (schematic, not to scale).
The ellipsis is cofinal in \(\varepsilon_0\); it
does not assert a uniform finite bound on exponent-tower height.}
\label{fig:p06-ordinals-epsilon-zero}
\end{figure}

% END HISTORICAL SLICE 288-460

\section{Further rank consequences and hierarchy diagram}

% BEGIN HISTORICAL SLICE 505-514
\begin{hbtheorem}{thm:p06-rank-recursion}{cited=Jech2003}{Rank recursion and membership decrease}
For every set \(x\),
\[
 \operatorname{rk}(x)=
 \sup\{\operatorname{rk}(y)+1:y\in x\}.
\]
Consequently \(y\in x\) implies
\(\operatorname{rk}(y)<\operatorname{rk}(x)\), and
\(\operatorname{rk}(\alpha)=\alpha\) for every ordinal \(\alpha\).
\end{hbtheorem}
% END HISTORICAL SLICE 505-514

% BEGIN HISTORICAL SLICE 530-564
\begin{figure}[htbp]
\centering
\begin{tikzpicture}[
  level/.style={draw,rounded corners,fill=black!4,minimum width=4.4cm,
                minimum height=0.55cm,align=center,font=\small},
  note/.style={align=left,font=\scriptsize}
]
  \draw[-{Latex[length=2.5mm]},thick] (0,0) -- (0,7.35)
    node[above,font=\small]{ordinal index};
  \node[level] (v0) at (3.0,0.35) {$V_0=\varnothing$};
  \node[level] (v1) at (3.0,1.15) {$V_1=\mathcal P(V_0)$};
  \node[level] (v2) at (3.0,1.95) {$V_2=\mathcal P(V_1)$};
  \node[note] at (3.0,2.65) {$\vdots$};
  \node[level] (vw) at (3.0,3.35)
    {$V_\omega=\bigcup_{n<\omega}V_n$};
  \node[level] (vws) at (3.0,4.15) {$V_{\omega+1}=\mathcal P(V_\omega)$};
  \node[note] at (3.0,4.85) {$\vdots$};
  \node[level] (vk) at (3.0,5.55)
    {$V_\kappa=\bigcup_{\alpha<\kappa}V_\alpha$};
  \node[level] (vks) at (3.0,6.35) {$V_{\kappa+1}=\mathcal P(V_\kappa)$};
  \foreach \y/\lab in {0.35/$0$,1.15/$1$,1.95/$2$,3.35/$\omega$,
                         4.15/$\omega+1$,5.55/$\kappa$,6.35/$\kappa+1$}
    {\draw (-0.08,\y)--(0.08,\y) node[left=3pt]{\lab};}
  \node[note,right=7pt of vk.east,text width=4.0cm] {
    Conditional large-cardinal marker: if $\kappa$ is strongly inaccessible,
    then $(V_\kappa,\in)\entails\mathsf{ZFC}$.  A measurable $\kappa$ is, in
    particular, inaccessible.\\[4pt]
    Since $\operatorname{rk}(\kappa)=\kappa$, the ordinal $\kappa$ itself
    first occurs as an element of $V_{\kappa+1}$.};
\end{tikzpicture}
\caption{The cumulative hierarchy and its rank axis.  The
large-cardinal annotation is conditional: the figure locates such a cardinal
if one exists and asserts no large-cardinal existence theorem.}
\label{fig:p06-cumulative-hierarchy}
\end{figure}
% END HISTORICAL SLICE 530-564
